\documentclass[a4paper, oneside, 12pt]{book}
\usepackage{graphicx}
\usepackage{float}
\usepackage{adjustbox}
\usepackage{listings}
\usepackage{fullpage}
\usepackage{booktabs}
\usepackage{subfig}
\usepackage{enumerate}
\usepackage{pgfplots}
\pgfplotsset{compat=1.8}
\usepackage{tabularx}
\usepackage{multirow}
\usepackage[T1]{fontenc}
\usepackage{inputenc}
\usepackage{hyperref}
\usepackage[italian]{babel}
\usepackage{longtable}
\usepackage{mathtools}
\usepackage{amssymb,amsmath,amsthm}

\begin{document}

\begin{titlepage}
\begin{center}

\includegraphics[scale=0.6]{Immagini/logo_dieti.png}
\vspace*{1cm}

{\Huge{\textbf{Progettazione e sviluppo di un applicativo in Java per la gestione di conferenze scientifiche}}}

\vspace{1cm}

{\Large Progetto d'esame per il corso di Object Orientation}


\vfill

{\Large Caporaso Antonio} \qquad \Large{Di Fusco Giorgio}\\
\texttt{matricola} \qquad \texttt{N86004389}

\vspace{1cm}

\includegraphics[scale=0.6]{Immagini/LogoUnina.png}

\vspace{1cm}
\begin{small}
Corso di Laurea Triennale in Informatica\\
Dipartimento di Ingegneria Elettrica e delle tecnologie dell'Informazione\\
Università degli Studi di Napoli Federico II\\
\textsc{A.A 2022/2023}
\end{small}
\end{center}
\end{titlepage}
\newpage
\textit{Questa pagina è stata lasciata intenzionalmente vuota.}
\newpage
\tableofcontents
\listoffigures
\listoftables
\end{document}
